% $REPORTS = /Volumes/Crucial X10/mootcourt-data/reports. Ablation reports: $REPORTS/ablations/backtest-upworthy-20261103-4bfb354b87e3--<setting>.json, each with tables.aggregate[cell=ALL] rows for P1 and P2.
\section{Ablations and Learning Curve}
\label{sec:ablations}

\paragraph{Does any setting of the precedent rung beat the prior?}
Table~\ref{tab:ablations-upworthy} (Appendix~\ref{app:tables}) changes one setting of P2 at a time on the same nine origins and 13,247 rows.
% src: $REPORTS/ablations/...--lambda-0.0.json P2 logscore -0.8812 brier_skill 0.0015; --lambda-0.2.json -0.9662, -0.0558; --lambda-0.5.json -0.9725, -0.0614; --lambda-1.0.json -0.9922, -0.0787
Fixing the precedent weight makes the loss monotonic: $\lambda$ of 0, 0.2, 0.5 and 1 give logarithmic scores of minus 0.881, minus 0.966, minus 0.973 and minus 0.992 and Brier skills of 0.001, minus 0.056, minus 0.061 and minus 0.079; $\lambda = 0$ reproduces P1 to the fourth decimal, as the construction requires.
% src: --bandwidth-mult-0.125.json P2 logscore -0.8892 brier_skill -0.0012 (-0.0048, 0.0026); --bandwidth-mult-0.25.json -0.8956; --bandwidth-mult-0.5.json -0.8961; --bandwidth-mult-1.0.json -0.8963; main run -0.8957
Fixing the bandwidth multiplier at 1/8, 1/4, 1/2 or 1 gives minus 0.889, minus 0.896, minus 0.896 and minus 0.896; the narrowest kernel is the least bad (Brier skill minus 0.001, minus 0.005 to 0.003), still does not beat the prior, and is better than the selected bandwidth by about 0.007 nats, so the selection also overfits the bandwidth.
% src: --features-text.json P2 logscore -0.8961 brier_skill -0.0079; --features-structured.json -0.8965, -0.0100; main run (both) -0.8957, -0.0078
Text features alone, structured features alone and both are within 0.001 nats of each other (minus 0.8961, minus 0.8965 and minus 0.8957): the headline text adds nothing to the structured context, and neither adds anything to the prior.
% src: --p2-selection-forward.json P2 logscore -0.8909 brier_skill -0.0001 (-0.0044, 0.0037) direction_accuracy 0.5759; --p2-selection-forward+recency-select.json -0.8910, -0.0002, p2_selection_by_cell half_life_days_median 90, 180, 90; --recency-select.json -0.8957, -0.0058, half-lives 90, 135, 90
Forward-chaining selection, which validates a candidate only on rows that start after an inner cut and so never lets a validation row borrow a precedent that did not yet exist, is the most principled rule and the least bad variant: minus 0.891 and a Brier skill of minus 0.000 (minus 0.004 to 0.004), with or without a recency kernel (half-lives of 90 to 180 days where it uses one; under leave-one-test-out recency moves the score by less than 0.001 nats). It ties the prior and does not beat it, so the default stays the rule the plan named; the rule for changing it (an advantage whose 90 percent interval clears zero) is in the method notes.
% src: --p2-guard-1.0.json P2 logscore -0.9293 brier_skill -0.0251; --p2-guard-off.json -0.9360, -0.0311; main run (guard 2.0) -0.8957, -0.0078
% src: PENDING reports/paper-extras-<snapshot>.json guard_admission_rates (the 39 and 8 percent on no-signal synthetic cells; today docs/method.md call 71)
Loosening the guard to one standard error gives minus 0.929 and a Brier skill of minus 0.025; removing it, minus 0.936 and minus 0.031; the two-standard-error guard of the main run, minus 0.896 and minus 0.008. On synthetic ledgers with no precedent signal the one-standard-error guard admitted a positive weight in 39 percent of cell fits and the two-standard-error guard in 8 percent, which is why the default is two.

\paragraph{Settings of the prior.}
% src: --kappa-0.0.json P1 logscore -0.8815; --kappa-5.0.json -0.8802; --kappa-80.0.json -0.8788; main run (kappa 20) -0.8812; --prior-normal.json P1 logscore -0.9572 brier_skill -0.0158 (-0.0215, -0.0105)
The pooling strength barely matters: $\kappa$ of 0, 5, 20 and 80 give P1 logarithmic scores of minus 0.882, minus 0.880, minus 0.881 and minus 0.879. The prior's shape matters: a normal prior in place of the Student-t costs 0.076 nats (minus 0.957) and turns the prior's Brier skill negative (minus 0.0158, minus 0.0215 to minus 0.0105), so the heavy tails of headline effects are real.
% src: $REPORTS/ablations/backtest-asos-20261103-6bb5a15c5ff0--prior-normal.json P1 logscore 4.3318 against main 4.6095, brier_skill -0.3678 (-0.6886, -0.1409); --kappa-0.0.json 4.5891, --kappa-5.0.json 4.6016, --kappa-80.0.json 4.6168
The same two findings hold on the deployment-shaped archive (Table~\ref{tab:ablations-asos} in Appendix~\ref{app:tables}): the normal prior loses 0.278 nats and falls to a Brier skill of minus 0.368, and $\kappa$ moves the score by under 0.03 nats.


\paragraph{Learning curve.}
% src: $REPORTS/learning-curve-upworthy-20261103-4bfb354b87e3.json final_origin 2015-01-01, heldout_rows 1143, heldout_tests 318, fractions [0.05, 0.1, 0.25, 0.5, 1.0], seeds 3; rows rung=P1 n_fit_rows 601 (fraction 0.05) and 11997 (1.00), skill 0.0266 (0.0131, 0.0429) and 0.0295 (0.0141, 0.0492)
Table~\ref{tab:learning-upworthy} (Appendix~\ref{app:tables}) subsamples the history visible at the final origin (January 2015) by whole tests at five fractions, three seeds each, and scores the same 1,143 held-out rows in 318 tests every time. P1 sits at a Brier skill of about 0.03 at every history size, 0.027 (0.013 to 0.043) with 601 history rows and 0.030 (0.014 to 0.049) with 11,997: twenty times more history adds nothing measurable, as a prior that has learned a base rate and a spread should behave.
% src: same file, rows rung=GBM skill -0.2836 (fraction 0.05) to -0.1161 (1.00); rung=local_eb 0.0090 (-0.0029, 0.0230) to 0.0025 (-0.0220, 0.0320); rung=text_lr 0.0004 (-0.0114, 0.0103) to -0.0100 (-0.0366, 0.0126); rung=P2 equals rung=P1 at every fraction
The gradient-boosted challenger improves with history (minus 0.284 to minus 0.116) and never reaches zero; the local empirical-Bayes set and the text classifier stay at zero within their intervals at every fraction; and P2 repeats P1 because the guard holds $\lambda$ at zero at every history size.
% src: $REPORTS/gates.json by_target 0.90 n_400 375; learning-curve-upworthy rows n_fit_rows minimum 601
This quarter's skill of 0.03 clears zero while the pooled skill over nine origins (Table~\ref{tab:main-upworthy}) does not; the two do not disagree, because the pooled figure averages nine quarters whose fits saw less history and whose base rates differ, while the learning curve scores one quarter with the fullest history and its bootstrap over 318 tests does not include the noise of choosing that quarter. The smallest history (601 rows) is above the 375-row gate, so no fraction here is undecidable.
% src: $REPORTS/learning-curve-asos-20261103-6bb5a15c5ff0.json heldout_rows 56, heldout_tests 9, skipped_runs {0.05: 3}; rows rung=P1 skill -0.0736 (-0.1805, 0.0088) at fraction 0.10 and -0.1824 (-0.5001, -0.0368) at 1.00
On the deployment-shaped archive the final origin holds 56 rows in 9 tests, the five-percent subsample cannot be fit, and every interval is wide (P1 from minus 0.074, minus 0.180 to 0.009, to minus 0.182, minus 0.500 to minus 0.037); the curve is produced for completeness (Table~\ref{tab:learning-asos}, Appendix~\ref{app:tables}) and decides nothing.

